% !TeX root = ../main.tex
\begin{center}
    \textbf{\Large ABSTRACT}
\end{center}

\vspace{0.5cm}

\noindent\textbf{Thesis title:} An LMS for Automatically Generating
Micro-Content \& Quizzes Using Generative AI

\smallskip

\noindent\begin{tabular}{@{}p{2.4cm}@{}l@{}}
  \textbf{Advisors:} & Dr. Nguyễn Quang Hùng \\
                     & Assoc. Prof. Dr. Thoại Nam \\[4pt]
  \textbf{Students:} & Phùng Xương Cận -- 2210348 \\
                     & Nguyễn Thanh Bình -- 2252082 \\
                     & Trần Ngọc Thiện -- 2134140 \\
                     & Nguyễn Ngọc Thành Đạt -- 2111013 \\
\end{tabular}

\vspace{0.7cm}

Converting textbooks, slides, and academic materials into concise learning
content, assessment questions, and review resources often requires significant
instructor effort. Popular AI tools can summarize content or generate
questions, but they often lack academic review workflows, LMS integration, and
mechanisms for aligning questions with course learning outcomes.

This thesis proposes, designs, and prototypes a Generative-AI-extended LMS
system that allows instructors to ingest materials, automatically analyze and
chunk them, generate Micro-Content and quizzes, and then review the generated
content before publishing it to learners. The system follows Hexagonal
Architecture/DDD and is organized as a multi-runtime stack consisting of a
Next.js frontend/BFF, NestJS and FastAPI backend services, asynchronous
workers, PostgreSQL, MinIO, Qdrant Vector Database, and Neo4j Knowledge
Graph. The Curriculum-Aware layer is designed to analyze the course syllabus,
extract learning outcomes, and link content chunks to Learning Outcomes in
order to support learning-objective-aligned question generation.

The core AI components include document parsing, heading-based chunking,
multilingual embeddings, vector retrieval, RAG, design-oriented GraphRAG, and
Bloom-taxonomy-aware quiz generation. The system targets LMS integration
through LTI 1.3/OIDC; Canvas LMS is used as the primary learning platform,
while the AI system operates as an independent External Tool. Because the
system follows the IMS LTI standard, the architecture is technically
compatible with other LTI 1.3-enabled LMS platforms, although that extension is
outside the scope of this thesis. The evaluation chapter defines functional,
Canvas/LTI integration, retrieval-quality, generated-content-quality,
performance, and limitations-evaluation criteria. Metrics that have not yet
been measured experimentally are explicitly marked for collection during the
pilot phase.

\vspace{0.5cm}

% \noindent\textbf{\large Keywords:} Micro-learning, Large Language Model,
% Retrieval-Augmented Generation, Knowledge Graph, Bloom Taxonomy, Vector
% Database, LTI, Learning Management System.
